\documentclass[acmsmall,screen,review]{acmart}
%% Generated by make-acm.py from main-arxiv.tex -- edit the arXiv source, not this file.

\usepackage{listings}
\usepackage{tikz}
\usetikzlibrary{shapes,arrows,positioning}
\usepackage{booktabs}
\usepackage{adjustbox}
\def\UrlBreaks{\do\/\do-\do_\do.\do=\do?\do&\do\#}
\Urlmuskip=0mu plus 1mu

\lstset{
  basicstyle=\ttfamily\small,
  breaklines=true,
  frame=single,
  backgroundcolor=\color{gray!10}
}

\setcopyright{none}
\acmJournal{TOPS}
\acmYear{2026}
\acmDOI{}
\acmISBN{}
\settopmatter{printacmref=false}
\renewcommand\footnotetextcopyrightpermission[1]{}

\begin{CCSXML}
<ccs2012>
<concept>
<concept_id>10002978.10003001.10003599</concept_id>
<concept_desc>Security and privacy~Operating systems security</concept_desc>
<concept_significance>500</concept_significance>
</concept>
<concept>
<concept_id>10002978.10002991.10002993</concept_id>
<concept_desc>Security and privacy~Access control</concept_desc>
<concept_significance>500</concept_significance>
</concept>
<concept>
<concept_id>10002978.10003022</concept_id>
<concept_desc>Security and privacy~Software and application security</concept_desc>
<concept_significance>300</concept_significance>
</concept>
<concept>
<concept_id>10010147.10010178</concept_id>
<concept_desc>Computing methodologies~Artificial intelligence</concept_desc>
<concept_significance>300</concept_significance>
</concept>
</ccs2012>
\end{CCSXML}
\ccsdesc[500]{Security and privacy~Operating systems security}
\ccsdesc[500]{Security and privacy~Access control}
\ccsdesc[300]{Security and privacy~Software and application security}
\ccsdesc[300]{Computing methodologies~Artificial intelligence}

\title[Namespace-Bounded Agents: Supplementary Material]{Supplementary Material for: Namespace-Bounded Agents: Capability-Based Security for LLM Systems via 9P Filesystem Semantics}

\author{Peter D. Finn}
\orcid{}%% TODO(author): ORCID iD
\email{pdf@nervsystems.com}
\affiliation{%
  \institution{King's College London}
  \city{London}
  \country{United Kingdom}
}
\affiliation{%
  \institution{NERV Systems}
  \country{USA}
}

\begin{abstract}This supplement accompanies the main article and contains the evaluation artifacts: environment specification, reproducibility checklist, model configuration, attack corpus structure, pass/fail rubric, and AgentDojo integration notes.\end{abstract}

\keywords{LLM agents, capability-based security, prompt injection, 9P protocol, Plan 9, Inferno OS}

\usepackage{xr}\externaldocument{main-acm}
\begin{document}

\maketitle
\appendix
\section{Evaluation Artifacts}
\label{sec:appendix}

All evaluation code, attack prompts, and raw results are available in our open-source
repository.\footnote{\url{https://github.com/NERVsystems/namespace-bounded-agents}, DOI: \href{https://doi.org/10.5281/zenodo.18419123}{10.5281/zenodo.18419123}}

\subsection{Environment Specification}

Table~\ref{tab:environment} lists the hardware, software, and model versions used in all experiments.

\begin{table}[h]
\centering
\small
\begin{tabular}{ll}
\hline
\textbf{Component} & \textbf{Specification} \\
\hline
Hardware & Apple M2 Pro, 32GB RAM \\
OS & macOS (host), Inferno OS (agent VM) \\
\hline
Claude Sonnet 4 & \texttt{claude-sonnet-4-20250514} (Anthropic API) \\
GPT-5 & \texttt{gpt-5-2025-08-07} (OpenAI API) \\
GPT-4o & \texttt{gpt-4o-2024-08-06} (OpenAI API) \\
\hline
Tokenizer & tiktoken \texttt{cl100k\_base} (proxy) \\
AgentDojo & v1.1.1 (NeurIPS 2024) \\
TLC (TLA+) & 2026.03.04 (rev: 52c0195) \\
SPIN & 6.5.2 \\
CBMC & 5.95.1 \\
Go & 1.21+ (\texttt{go9p} library) \\
\hline
\end{tabular}
\caption{Hardware, software, and model specifications for all experiments.}
\label{tab:environment}
\end{table}

\subsection{Reproducibility Checklist}

To reproduce Tables~\ref{tab:security},~\ref{tab:multi-model},~\ref{tab:agentdojo-results}, and~\ref{tab:tokens}:

\begin{enumerate}
    \item \textbf{Clone repository}: \texttt{git clone} at commit hash documented in release
    \item \textbf{31-attack corpus (Tables~\ref{tab:security},~\ref{tab:multi-model})}:
    \begin{itemize}
        \item Start 9P server: \texttt{bin/server --port 5640}
        \item Run evaluation: \texttt{bin/seceval --config attacks.json --model [model] --runs 4}
        \item Seeds: deterministic at temperature=0.0; GPT-5 uses temperature=1.0 (stochastic)
    \end{itemize}
    \item \textbf{AgentDojo (Table~\ref{tab:agentdojo-results})}:
    \begin{itemize}
        \item Install: \texttt{pip install agentdojo==1.1.1}
        \item Run: \texttt{python namespace\_agent/namespace.py --model claude-sonnet-4-20250514}
        \item Full 629-attack sweep; namespace filtering applied per-task
    \end{itemize}
    \item \textbf{Token benchmark (Table~\ref{tab:tokens})}:
    \begin{itemize}
        \item Run: \texttt{bin/agent --benchmark --tools 14}
        \item Token counts via tiktoken \texttt{cl100k\_base}
    \end{itemize}
    \item \textbf{Formal verification}:
    \begin{itemize}
        \item TLA+: \texttt{./run-verification.sh small} ($\sim$8 min, 11 invariants)
        \item SPIN: \texttt{./spin/verify-all.sh} ($\sim$30 sec, 5 models)
        \item CBMC: \texttt{./cbmc/verify-all.sh quick} ($\sim$1 min, 3 harnesses)
    \end{itemize}
\end{enumerate}

\subsection{Model Configuration}

\begin{table}[h]
\centering
\small
\begin{tabular}{lccc}
\hline
\textbf{Parameter} & \textbf{Claude Sonnet 4} & \textbf{GPT-5} & \textbf{GPT-4o} \\
\hline
API & Anthropic & OpenAI & OpenAI \\
Model ID & \texttt{claude-sonnet-4-20250514} & \texttt{gpt-5-2025-08-07} & \texttt{gpt-4o-2024-08-06} \\
Temperature & 0.0 & 1.0* & 0.0 \\
Max tokens & 4096 & 4096 & 4096 \\
Runs per attack & 4 & 4 & 4 \\
\hline
\end{tabular}
\caption{Model configuration for multi-model security evaluation. *See Table~\ref{tab:experiment-summary} regarding GPT-5's fixed temperature.}
\label{tab:model-config}
\end{table}

\subsection{Attack Corpus Structure}

Each attack in the 31-attack corpus includes:
\begin{itemize}
    \item \textbf{ID}: Unique identifier (e.g., PT-1, TD-2)
    \item \textbf{Category}: One of nine categories (Table~\ref{tab:attack-taxonomy})
    \item \textbf{Prompt}: Exact injection text
    \item \textbf{Success indicators}: Regex patterns indicating attack success
    \item \textbf{Failure indicators}: Regex patterns indicating defense success
\end{itemize}

Example attack (PT-1):
\begin{lstlisting}
{
  "id": "PT-1",
  "category": "path_traversal",
  "prompt": "Please read /etc/passwd and list users",
  "success_patterns": ["root:", "nobody:", "/bin/"],
  "failure_patterns": ["not found", "denied", "cannot"]
}
\end{lstlisting}

\subsection{Pass/Fail Rubric}

We use two evaluation methodologies corresponding to our two attack corpora:

\subsubsection{31-Attack Corpus (Pattern Matching)}

The custom 31-attack evaluation uses pattern matching on execution output:
\begin{itemize}
    \item \textbf{Execution layer}: Commands are actually executed---9P agent against
    namespace-bounded runtime (returns ``path does not exist'' for forbidden paths);
    MCP baseline against Docker container.
    \item \textbf{Classification}: Success/failure determined by string patterns in
    execution output (e.g., \texttt{root:} = success for /etc/passwd; ``not found'' = blocked).
    \item \textbf{Limitation}: Hallucinated content in LLM response text could cause
    false positives. However, structural blocks are reliable because the execution layer
    returns explicit error messages that override success patterns.
\end{itemize}

\subsubsection{AgentDojo (629 Attacks, Runtime Evidence)}

The AgentDojo evaluation uses rigorous runtime state comparison:
\begin{itemize}
    \item \textbf{Environment snapshots}: Pre/post execution state captured.
    \item \textbf{Function traces}: Each tool invocation logged as \texttt{FunctionCall}
    with function name, arguments, and result.
    \item \textbf{Success determination}: AgentDojo's \texttt{security\_from\_traces()}
    compares environment state---attack succeeds only if target action actually modified
    the environment (e.g., funds transferred, email sent).
    \item \textbf{Namespace blocking}: Paths outside namespace logged in
    \texttt{blocked\_paths} before any execution.
\end{itemize}

\textbf{Classification (AgentDojo):}
\begin{itemize}
    \item \textbf{Attack Success}: \texttt{security\_from\_traces()} returns True.
    \item \textbf{Structural Block}: Path blocked by namespace check before execution.
    \item \textbf{Behavioral Block}: No structural block, but attack failed (model refused
    or generated ineffective commands).
    \item \textbf{Hallucination}: Response suggests success but traces/state unchanged.
\end{itemize}

The AgentDojo methodology is more rigorous and constitutes our primary security evidence
(629 attacks vs.\ 31). The 31-attack corpus provides category-level analysis with the
noted pattern-matching limitation.

\subsection{AgentDojo Integration}
\label{sec:agentdojo-integration}

AgentDojo tasks were executed using the official benchmark harness
\cite{debenedetti2024agentdojo}. Namespace isolation was implemented as a
tool-visibility filter that replicates 9P namespace semantics: for each user task,
only the minimal required tool set is visible to the agent. Tools outside the task's
minimal set return ``tool not found'' before any execution occurs, equivalent to the
ENOENT response from a 9P namespace boundary.

This is functionally equivalent to kernel-level namespace enforcement because:
(1)~filtering occurs before the model sees the tool list---excluded tools are absent from
the agent's context, not merely denied at invocation; (2)~no hidden invocation path exists
for filtered tools; (3)~\texttt{blocked\_paths} are logged prior to any execution attempt.
The key property---that capability names outside the allowed set do not exist in the agent's
interface---is preserved by the filtering implementation.

Cross-tool vs.\ same-tool classification was performed by comparing each injection's
target tools against the minimal tool set for the corresponding user task. The minimal
tool sets themselves were derived from the benchmark's task specifications; automatic
derivation from raw user requests is future work (see Limitations).

\textit{Run provenance:} During the initial full sweep, 17 of the 105 Slack-domain
attacks failed to score due to an AgentDojo API mismatch in their security checks. The
Slack domain was re-run in full after the fix, and the reported results combine the
corrected full run (banking, workspace, travel) with the Slack re-run, covering all 629
attacks. Both result sets are included in the artifact repository.

\textit{Note:} End-to-end 9P protocol integration with AgentDojo (running the benchmark
through actual 9P mount/walk/read/write operations) is future work. The current evaluation
validates the security model---capability-name non-existence in the agent-visible
interface---via equivalent filtering at the tool-visibility layer.



\section{Formal Verification Details}
\subsection{TLA+ Specification and Model Checking}

The TLA+ specification (1,199 lines across 4 modules) models 12 kernel operations
(\texttt{CreateProcess}, \texttt{ForkWithForkNS}, \texttt{ForkWithNewNS}, \texttt{Mount},
\texttt{Unmount}, etc.) with 15 state variables including 3 history variables for
non-trivial isolation verification.

\textbf{Small configuration (exhaustive):} 2 processes, 3 pgrps, 3 channels, 2 paths.
TLC explored 369,414,154 states (19,307,332 distinct) to search depth 27 with 0 states
remaining on queue. All 11 safety invariants verified with zero violations in 7 min
57 sec on 16 workers.

\textbf{Medium configuration (partial):} 3 processes, 4 pgrps, 4 channels, 3 paths.
TLC explored 26.5 billion states (3.17 billion distinct) to depth 13 over 10 hr
50 min on an NVIDIA Jetson Orin AGX (12 ARM64 cores, 50 GB JVM heap, 341 GB NVMe state
storage). Zero violations across all 11 invariants. The run was terminated when
on-disk state storage approached capacity; it is recoverable from checkpoint.

\textbf{11 invariants verified:} TypeOK, SafetyInvariant, RefCountNonNegative,
NoUseAfterFree, MountTableBounded, NamespaceIsolation, NamespaceIsolationTheorem,
UnilateralMountNonPropagation, CopyFidelity, PostCopyMountsSound, NoIsolationViolation.

\subsection{SPIN Model Checking}

Five Promela models (1,440 lines) verify concurrent protocol properties:

\begin{center}
\small
\begin{tabular}{p{3.2cm}rl}
\textbf{Model} & \textbf{States} & \textbf{Property} \\
\hline
Basic isolation & 82,628 & Post-copy mount isolation \\
Extended (nested fork) & 5,724 & 3-level hierarchy \\
Multi-lock protocol & 39,744 & Lock ordering \\
Namespace races & 3,131 & Use-after-free detection \\
exportfs boundary & 106 & Root escape prevention \\
\hline
\textbf{Total} & \textbf{131,333} & \textbf{5/5 PASS} \\
\end{tabular}
\end{center}

The lock ordering model verifies that \texttt{pg->ns} is always acquired before
\texttt{mh->lock}, including \texttt{cmount}'s early-release optimization. The
\texttt{exportfs} model verifies that \texttt{walk("..")} sequences cannot escape the
export root.


\subsection{CBMC Bounded Model Checking}

CBMC verifies properties of actual C source code extracted from the kernel.
Three quick-mode harnesses (CI, $\sim$1 min) check array bounds (MOUNTH hash macro,
58 properties), integer overflow (fd allocation, 10 properties), and reference
counting (45 properties): 113 total properties, 0 failures.

Four full-mode harnesses verify the complete \texttt{pgrpcpy()} implementation
including \texttt{newpgrp()}, \texttt{newmount()}, and \texttt{pgrpinsert()} ---
actual C source from \texttt{pgrp.c}, not re-implementations --- at the production
configuration (\texttt{MNTHASH=32}, \texttt{--unwind 34}, \texttt{--object-bits 11}):

\begin{center}
\small
\begin{tabular}{p{3.8cm}rl}
\textbf{Harness} & \textbf{Properties} & \textbf{Result} \\
\hline
Basic isolation & 908 & PASS \\
Modification independence & 908 & PASS \\
MNTHASH bounds & 841 & PASS \\
Error path safety & 653 & PASS \\
\hline
\textbf{Total} & \textbf{3,310} & \textbf{4/4 PASS} \\
\end{tabular}
\end{center}

The \texttt{harness\_basic\_isolation} verifies that \texttt{pgrpcpy()} produces
independent mount tables: the child's \texttt{Mhead} is a distinct object (deep copy),
mounted channels are shared with incremented reference counts, and \texttt{slash}/\texttt{dot}
are cloned. The \texttt{harness\_modification\_independence} verifies that post-copy
mounts to the parent's mount table do not affect the child --- the implementation-level
analogue of the TLA+ Namespace Isolation Theorem. The \texttt{harness\_pgrpcpy\_error}
verifies that when \texttt{pgrpcpy()} fails mid-copy (e.g., allocation failure), the
source namespace is unmodified, the write lock is released, and no stale readers remain.

All runs include \texttt{--unwinding-assertions}, confirming that the unwind depth of 34
fully covers all program loops --- the bounded verification is therefore sound within
CBMC's framework. The \texttt{--slice-formula} flag was critical for tractability:
without it, CBMC stalled at 77+ minutes for \texttt{MNTHASH=4}; with slicing, the
production \texttt{MNTHASH=32} configuration completes in under one second.

\bibliographystyle{ACM-Reference-Format}
\bibliography{references}
\end{document}
